\documentclass{article}
\usepackage{amsmath}
\begin{document}
\title{A removed Gazebo entity's numeric identity can recur across a world reset}
\maketitle

\paragraph{Result and scope.}
This selected-source replay tests Gazebo Sim 9 pinned to
\texttt{62f198e0478159549b748c4b4c66f0cbe2fcbf6f}. The real allocator,
single-remove processing, private copy, vertex diff, diff application
and reset methods run in a serial graph fixture. Vertices have no
component data, view contents or joints in this test; single-entity
removal marking is supplied by the fixture for that scope. The pinned
runner is inspected to establish that it stores the initial entity
manager and calls this reset before stepping. The test does not
instantiate a full server, engine, plugin restart, or remote transport.

\paragraph{Exact allocator cut.}
Within this scope, removing an entity leaves the entity counter unchanged.
However, reset first copies the initial manager, including its counter
$b$. It compares the present and initial graph by vertex identifier, then
adds current-only vertices to the temporary graph and marks them for
removal. Adding these vertices raises the counter up to their maximum
identifier. Vertices already absent before the reset do not take this path.

Let $S$ be the identifiers present just before the reset that were added
after the initial snapshot, all above $b$. Assume no counter override,
distributed allocation, reused identifiers or overflow before the test.
All initial vertex identifiers are at most $b$. The counter after reset
and normal removal processing is exactly
\[
c_{\mathrm{reset}}=\max\bigl(\{b\}\cup S\bigr).
\]
Initial-only vertices are already in the copy; remaining initial vertices
do not raise the count. No component or view behavior is needed for this
claim. The next ordinary allocations therefore repeat the old suffix
of identifiers above $c_{\mathrm{reset}}$ that were previously issued.
A retained high vertex prevents that particular suffix from repeating,
even if reset subsequently removes the high vertex. This boundary agrees
with all 2,550 source cases covering every live subset through eight
new entities at five initial counts, including 1,275 cases with a
reused suffix.

\paragraph{Why it matters for a receiver.}
The smallest exercised witness starts at counter 10, creates entity 11,
fully removes it, resets and allocates entity 11 again. The number can
now name a new actor incarnation, distinct from the subject of an old
actual-effect receipt. In the explicit receiver negative control, an
external archive retains a receipt from before reset. A replacement
lookup by world, entity number and reused job number admits the old
receipt for work the new actor has not done. A stale receipt and a
legitimate new receipt can have the same source-visible tuple in that
interface. This is an authored archive/admission example, not a fault
in a stock receipt API or a claim that transport buffers survive reset.

One adequate repair for the stated model is to bind admission to the
authoritative world incarnation as well as to the body and operation:
the incarnation counter or name must be outside the snapshot that can
rewind. If the authority cannot establish the current incarnation,
it cannot use an id-only receipt to admit this work. Operations that are
intentionally carried across reset require their own explicit mapping;
the example does not make all cross-reset evidence invalid. Other valid
designs include a non-reusable operation identifier or a barrier proving
that no earlier receipts remain available anywhere. Merely moving a
generation field into the rewindable checkpoint does not provide the
needed distinction. Both the explicit test and an equally informed direct
receiver agree.

\paragraph{Limits.}
This invalidates a potential blanket use of native numeric entity
identifiers as lifetime-stable appointment identifiers. It does not show
that every reset reuses identifiers, that the selected source case
physically repeats a move, or that an implementation already supplies a
durable incarnation authority. This is a scoped interface obstruction
and positive design requirement, not independent mission or architecture
acceptance.
\end{document}
